\section{Recovering identifiability by readout and metric design}
\label{sec:recovery}
Changing the metric or the readout changes the kernel, unlike the
equivalence-preserving transformations above. The design objective is
not to make every source symbol active, but to recover dependence
already visible in the centered features and to identify its limits.

\subsection{What can a metric change recover?}
At a fixed parameter point, write $H_v=D C(\theta)[v]$. By
Theorem~\ref{thm:gram}, the all-input kernel derivative vanishes exactly
when
\begin{equation}
 C^T M H_v+H_v^T M C=0.\label{eq:recoverability}
\end{equation}
This test separates a symmetry of a particular metric from a loss of
dependence in the features themselves.

\begin{proposition}[Irrecoverability over fixed metrics]
\label{prop:metricrecover}
Assume the finite-basis hypotheses of Theorem~\ref{thm:gram}. At a
point where $f_\theta$ is nonconstant, a direction $v$ is invisible to
every fixed positive-definite metric if and only if $H_v=0$,
equivalently $D_v f_\theta(x)$ is input-independent.
\end{proposition}
\begin{proof}
For an input pair set $a=f_\theta(x)-f_\theta(x')$ and
$b=D_v f_\theta(x)-D_v f_\theta(x')$. If $a^TMb=0$ on the open
positive-definite cone, linearity in $M$ implies
$ab^T+ba^T=0$. For $a\ne0$ this gives $b=0$. When $a=0$, compare
both inputs to a third input with a different feature value, which
exists by nonconstancy. Thus $b=0$ for every pair. Affine independence
of the data basis implies $H_v=0$. The converse is immediate.
\end{proof}
The nonconstant assumption is essential: at $C=0$, every
squared-distance first derivative vanishes regardless of $H_v$.
Also, the intersection of nullspaces over all metrics need not equal
the nullspace of any \emph{single} metric. The bound
$r\le A(A+1)/2$ continues to constrain each fixed metric.

At a fixed parameter point, the minors of Eq.~\eqref{eq:recoverability},
viewed as a map of $v$, are polynomials in the metric entries.
Consequently, if a rank is attainable in an open positive-definite
metric family, almost every metric in that family attains its maximum
rank. This applies separately within the diagonal family and to
finitely many parameter points simultaneously
(Appendix~\ref{app:recovery}). Random anisotropy is therefore a
necessary baseline; generic rank alone does not imply conditioning
or good finite-shot sensitivity.

\subsection{An exact boundary between metric and feature recovery}
The question here is whether a fixed metric on the existing features
can recover all four feature-visible directions, or whether the
representation itself must change.
For the four-qubit, one-layer Hubregtsen source circuit, the local-Z
features are $f_q(x)=-\sin\theta_q\cos x_{q\bmod2}$ for $q=0,\ldots,3$.
Its diagonal controlled-rotation tail changes none of these features
(the circuit derivation is given at Eq.~\eqref{eq:hubfeatures}).
The four retained angles vary independently; the remaining tail
angles do not enter these features. Assume that
$(\cos x_0,\cos x_1)$ ranges over a set with nonempty interior.
The following statement concerns generic differential rank, not
rank at every parameter point or globally unique parameter recovery.

\begin{proposition}[Fixed-quadratic ceiling and product recovery]
\label{prop:recoveryboundary}
For this four-qubit Hubregtsen/local-Z family with fixed
$\gamma>0$: (i) every fixed positive diagonal metric on the four
local means has generic kernel rank two; (ii) the maximum generic
rank over fixed dense positive-definite metrics on these means is
three, attained by Eq.~\eqref{eq:recoverymetric}; (iii) appending
$f_0f_2$ and $f_1f_3$ gives generic rank four for every fixed positive
diagonal metric on the resulting six features. These products equal
the corresponding ZZ means in this family and need no additional
ideal measurement information.
\end{proposition}
\begin{proof}
Write $a=\sin\theta_0$, $b=\sin\theta_1$,
$c=\sin\theta_2$, and $d=\sin\theta_3$. For diagonal weights, the coefficient
Gram on $h=(\cos x_0,\cos x_1)^T$ is
\[
 G=\diag(w_a a^2+w_c c^2,\ w_b b^2+w_d d^2).
\]
Its two entries have generically independent derivatives, proving (i).
For every dense metric the same Gram is symmetric $2\times2$, giving
rank at most three. With $u=(1,2,3,4)^T$,
\begin{equation}
 M=(32I+uu^T)/158\label{eq:recoverymetric}
\end{equation}
is positive definite with unit trace, and an exact nonzero rational
minor certifies rank three. Its value is
$-\frac{835727837184}{978271035584375}$ at the rational-phase point in
Appendix~\ref{app:recovery}, proving (ii).

The state before the diagonal tail is a product and the tail commutes
with every Z string. Consequently $f_{qr}(x)=f_q(x)f_r(x)$ exactly.
These ZZ features are a nonlinear lift of the original four means,
not new ideal information that was absent from those means.
For (iii), put $(u,v)=(\cos x_0,\cos x_1)$.
The data functions $1,u,v,u^2,v^2$ are independent. Four entries of
their coefficient Gram are
\begin{equation}
 \begin{aligned}
 \beta_u&=w_a a^2+w_c c^2,&\beta_{u^2}&=w_{ac}a^2c^2,\\
 \beta_v&=w_b b^2+w_d d^2,&\beta_{v^2}&=w_{bd}b^2d^2.
 \end{aligned}\label{eq:twoproductrecovery}
\end{equation}
The $(a,c)$ Jacobian block has determinant
$4acw_{ac}(w_a a^2-w_c c^2)$, and the $(b,d)$ block has the
analogous generically nonzero determinant. Multiplying by the four
cosine derivatives preserves generic nonvanishing, proving rank at
least four. Only the four pre-tail angles occur, giving the matching
upper bound.
\end{proof}
The experiment retained all six ZZ products and also
has exact minor certificates. Two products are sufficient; we do not
claim their minimality over all allowed feature/metric designs.

\begin{table}[t]
\caption{Exact recovery limits for the same four-qubit, one-layer
Hubregtsen source circuit. Ranks are generic, not statements at every
parameter point. The third row is the two-product analytic construction
of Proposition~\ref{prop:recoveryboundary}; the fourth is the
configuration used in the finite-shot experiment. All rows can use
one ZZZZ product basis. Both product lifts are also available from
exact local means by nonlinear classical processing.}
\label{tab:recoverymechanism}
\begin{tabular}{@{}>{\raggedright\arraybackslash}p{.74\columnwidth}r@{}}
\toprule
Features and fixed metric & Rank\\
\midrule
Four local Z, any positive diagonal metric & 2\\
Same four features, best dense metric & 3\\
Local Z plus two products, positive diagonal metric (analytic proof) & 4\\
Local Z plus six ZZ products, diagonal metric (experimental configuration) & 4\\
\bottomrule
\end{tabular}
\end{table}
\begin{corollary}[Fixed affine expansion cannot cross the ceiling]
\label{cor:affineceiling}
For the four local features of Proposition~\ref{prop:recoveryboundary},
let $\widetilde f=Lf+b$ have any finite output dimension, with $L,b$
fixed independently of inputs and trainable parameters. With any
fixed $\widetilde M\succ0$ and fixed bandwidth, its kernel has
generic differential rank at most three. Thus no such affine
expansion achieves rank four, whereas the two quadratic product
features in the proposition suffice.
\end{corollary}
\begin{proof}
The offset cancels and
\begin{equation}
 \Delta\widetilde f^T\widetilde M\Delta\widetilde f
 =\Delta f^T L^T\widetilde M L\Delta f.
\end{equation}
The effective metric may be singular, but the data basis is still
$(\cos x_0,\cos x_1)$. Its symmetric $2\times2$ coefficient Gram
has only three entries, bounding the differential rank by three.
\end{proof}
This is a restriction on fixed affine lifts, not on arbitrary
input-dependent processing or jointly trainable postprocessing.

At finite shots, direct ZZ parity averages and products of local-Z
sample means are different estimators of the same ideal features in
this family. We compare them on identical outcomes in
Sec.~\ref{sec:productcontrol}. General correlated states do not obey
this product identity. The result establishes a limit of fixed
quadratic metrics on the original features, not a necessity for
additional quantum measurement information. Equal basis count alone
does not imply equal estimation precision.

On ten previously studied source circuits from the YZ-CX,
Hubregtsen, and Kyriienko constructors, fixed diagonal local-XYZ
weights proportional to $1,\ldots,4$ on an evenly spaced grid attain
generic rank $K$, certified by exact Fourier-Gram minors. In
particular, the two- and four-qubit Kyriienko ranks change from
$4$ to $6$ and from $10$ to $12$. These are redesigned kernels on
existing circuits, not additional entries in the original 54-setting
census. The broader 30-setting metric survey is numerical; ten
single-basis designs also have exact full-rank certificates. Simple
random metrics and uniform-basis controls already attain many of the
same ranks, so these results do not establish a sophisticated selector's
superiority.

The next two sections give the circuit-specific certificates and
dimension census supporting this identifiability picture
(Secs.~\ref{sec:structure} and \ref{sec:census}). We then return to
our third question in Sec.~\ref{sec:shotdesign}: how well can the
retained or restored directions be resolved with finite shots?
Exact rank is an algebraic prerequisite, not a guarantee of statistical
resolution or predictive improvement.
